% ===========================================================================
% Appendix note: reconstruction-uncertainty bootstrap band methodology.
% Intended to REPLACE the \TODO in appendix_human.tex's "Reconstruction
% uncertainty" paragraph (app:human / app:mixture), and to be referenced by the
% _bands version of Figure~\ref{fig:smad-comparison} and by Figure~\ref{fig:ranking}.
% Generated by analysis/reconstruction_bands.py (seed 1299, B=2000).
% NOTE: fragment only; no file under writeup/ is edited by this task.
% ===========================================================================

\paragraph{Reconstruction-uncertainty bands (parametric bootstrap over the
calibration).}
Beyond the Monte-Carlo simulation noise conditional on a fixed calibration, we
represent uncertainty in the calibration itself by a parametric bootstrap over
the mixture parameters, exactly as committed in Section~\ref{sec:reconstruction}.
For each source--treatment pair we draw $B=2000$ vectors of
$(p_{\text{eq}}, p_{\text{over}}, p_{\text{under}}, \lambda, \sigma, \delta)$
from admissible ranges pinned to the reported statistics and their granularity,
propagate each draw through the reconstruction that defines the anchor, and
report percentile bands (both $2.5/97.5$ and $16/84$). The point anchors are
reproduced exactly at the central (published) calibration before any
perturbation, so the bands are centered on the numbers quoted in the main text.

The admissible ranges are: the direction weights
$(p_{\text{eq}}, p_{\text{over}}, p_{\text{under}})$ are re-drawn from the
sampling uncertainty of the reported over-/under-/at-value \emph{frequencies}
(a multinomial resample of the reported cell counts for
\citet{kagel1993independent}; a $\pm 5$ percentage-point range around each
reported behavioral rate for \citet{li2017obviously} and
\citet{breitmoser2022obviousness}, whose rates are stated only to that
granularity, e.g.\ ``${\sim}50\%$ dominant-strategy play'' admits $45$--$55\%$);
the mean bid-to-value ratio is re-drawn through the reported bid-regression slope
$\hat\beta$ within its Table~2 standard error ($\mathrm{SE}\approx 0.01$), which
is the dominant driver of the sealed-bid IPV bands (for FPSB the anchor is the
bid-shading gap $\lvert \bar r - \tfrac{N-1}{N}\rvert/\tfrac{N-1}{N}$, since human
bids sit essentially at value); the equilibrium noise $\sigma$ is re-drawn
through the reported $R^{2}$ (to which it is tied via
$\sigma^{2}=\hat\beta^{2}\mathrm{Var}(v)(1-R^{2})/R^{2}$) within its reporting
precision; the underbid bound $\delta$ is re-drawn over the parsimony range
$[0.10, 0.25]$; and for the clock/APV formats the reported mean absolute
deviation is re-drawn within its published standard error
($\mathrm{MAD}\sim\mathcal N(\mathrm{MAD}, \mathrm{SE}(\mathrm{MAD})^{2})$), which
reproduces the standard errors already carried in
Table~\ref{tab:mixture-calibration}'s source cells. The overbid scale $\lambda$
governs the shape rather than the level of the deviation mass and is pinned so
the central mixture reproduces the published anchor.

These are calibration bands, not sampling bands on human behavior: they quantify
how far the reconstructed anchor can move within what the source statistics leave
undetermined, and we continue to use the reconstructions only for orderings,
directions, and comparative statics---never levels or hypothesis tests. Under the
bands, the two load-bearing orderings are stable: the human first-price
$\gg$ second-price difficulty gap holds in $100\%$ of re-draws (SMAD ratio
$95\%$ band $[3.54, 5.65]$, centered on the $4.4\times$ point gap), and the
clock improvement over the sealed second-price benchmark holds in $100\%$ of
re-draws for the open clock (AC) and $92.6\%$ for the closed clock (AC-B), the
latter's wider band inherited from the large reported standard error on the
\citet{breitmoser2022obviousness} mean absolute deviation. Full per-anchor bands
are in \texttt{results/reconstruction\_bands/bands.csv} and
\texttt{bands\_summary.md}.
